\section{Algebra of local fields}
Local fields \(A(x)\in\mathcal{A}\) are operators in a QFT depending on only one space-time point. This set includes fundamental fields, but also their powers and derivatives. However, they are not freely generated: the algebra of local fields is subject to constraints (for example, equations of motion like \(\partial^2\phi = \phi^5\)). The set \(\mathcal{A}\) is called the algebra of local fields. Indeed:
\begin{itemize}
    \item Local fields add up and can be multiplied by scalars, so \(\mathcal{A}\) has the structure of a vector space. We can select a preferred basis \(\{\varphi_n(x)\}\) for this space, such that all fields can be expressed as a linear combination of this basis.
    \item Usually, we choose a basis of local fields that diagonalizes two commuting operators (since they commute, they can be simultaneously diagonalized):
          \begin{itemize}
              \item \textbf{Dilatation}: \(\tilde{D}\varphi_{n}(0)=\lambda^{d_{n}}\varphi_{n}(0)\). The eigenvalue \(d_{n}\) (or \(\Delta_n\)) is the scale dimension. This is the crucial number we are interested in, as it dictates how the field rescales and ultimately determines the critical exponents that characterize specific universality classes.
              \item \textbf{Spin}: \(S\varphi_{n}(0)=e^{is_{n}\theta}\varphi_{n}(0)\). Local fields can have various spins (\(s_n = 0, \tfrac{1}{2}, 1, \tfrac{3}{2}, \dots\)).
          \end{itemize}
    \item \(\varphi_{n}(x)\) is therefore uniquely characterized by this pair of numbers \(\{d_{n},s_{n}\}\).
    \item All other fields are linear combinations
          \[
              A(x)=\sum_{n=0}^{\infty}a_{n}\varphi_{n}(x), \quad \varphi_{0}\equiv\mathbb{I},
          \]
          where the equality here has to be understood in the \textit{weak sense}. This means that two operators are considered equal if they yield the exact same result when inserted into \textit{any} correlation function of the theory. You can have two formally different expressions, but if they do the same job on the full algebra of fields, they represent the same physical field. Formally, for any string of local fields \(X=A_{1}(x_{1})\dots A_{n}(x_{n})\in\mathcal{A}^{\otimes n}\), we have
          \[
              A(x)=B(x)\Leftrightarrow\langle A(x)X(x_{1},\dots,x_{n})\rangle=\langle B(x)X(x_{1},\dots,x_{n})\rangle.
          \]
\end{itemize}
2-point functions of local fields, at criticality, are fixed by roto-translation and scale invariance:
\[
    \mathcal{D}_{ij}(x_{1}-x_{2})=\langle\varphi_{i}(x_{1})\varphi_{j}(x_{2})\rangle=\frac{A_{ij}}{|x_{1}-x_{2}|^{d_{i}+d_{j}}},
\]
where \(A_{ij}=A\delta_{ij}\). So they are all zero if \(\phi_{j}\ne\phi_{i}\). The constant \(A\) can be fixed by normalizing the field \(\phi_{i}\):
\[
    \mathcal{D}_{ij}(x_{1}-x_{2})=\langle\varphi_{i}(x_{1})\varphi_{j}(x_{2})\rangle=\frac{\delta_{ij}}{|x_{1}-x_{2}|^{2d_{i}}}.
\]
Near criticality we have a mass scale \(m\sim\xi^{-1},\) i.e. the inverse of the correlation length. The expectation values of local fields (think e.g. to the order parameters) rescale as:
\[
    \langle\varphi_{i}(x)\rangle=\mu_{i}m^{d_{i}}.
\]
At criticality, i.e. in presence of scale invariance and \(\xi\rightarrow\infty\):
\[
    \langle\phi_{0}\rangle=\langle\mathbb{I}\rangle=1, \quad \langle\phi_{i}(x)\rangle=0 \quad (i=1,2,\dots).
\]

The vector space of local fields is raised to an algebra if a suitable notion of product between fields is established: the \textbf{Operator Product Expansion (OPE)}.
Any product of fields that are separated by some distance can be expanded as a sum of individual fields. Consider a composite (non-local) field \(A(x_{1})B(x_{2})\). In the limit \(x_{1}\rightarrow x_{2}\), we expect that it reproduces the properties of a local field \(C(x_{2})\), which can then be expanded in the basis of local fields. This is conceptually vital to understand objects like \(\phi^2\), which is rigorously defined as \(\phi(x_1)\phi(x_2)\) evaluated as \(x_1 \to x_2\). Note that \(x_1\) and \(x_2\) do not need to be infinitesimally close; the OPE holds as long as there are no other field singularities in the middle (i.e., within a certain radius of convergence).

Wilson assumed that for any pair of local fields \(A(x_{1})\) and \(B(x_{2})\) the OPE holds:
\[
    A(x_{1})B(x_{2})=\sum_{k=0}^{\infty}c_{k}(x_{1},x_{2})\varphi_{k}(x_{2}).
\]
In particular, for our chosen basis:
\[
    \varphi_{p}(x_{1})\varphi_{q}(x_{2})=\sum_{r}c_{pq}^{r}(x_{1},x_{2})\varphi_{r}(x_{2}).
\]
The functions \(c_{pq}^{r}(x_{1},x_{2})\) are called structure functions of the OPE algebra. Because the theory is \textit{translationally invariant}, these functions cannot depend on the absolute coordinates \(x_1\) and \(x_2\), but only on their difference \(x_1 - x_2\). Furthermore, \textit{rotational invariance} (Poincar\'e invariance in the Wick rotated Euclidean description) implies they depend only on the distance between the two points (the modulus, not the full vector):
\[
    c_{pq}^{r}(|x_{1}-x_{2}|)=c_{pq}^{r}(|x_{12}|).
\]
Scale invariance at criticality further restricts the function:
\[
    c_{pq}^{r}(x_{1},x_{2})=\frac{C_{pq}^{r}}{|x_{12}|^{d_{p}+d_{q}-d_{r}}}.
\]
The numbers \(C_{pq}^{r}\) are called structure constants.

Here and in the following we denote, for short,
\[
    x_{ij}=x_{i}-x_{j}.
\]
The set of all scaling dimensions \(d_n\) (the spectrum of the theory) together with the structure constants \(C_{pq}^{r}\) form the so-called \textbf{CFT data}. These quantities completely define the dynamical content of the theory. If they are known exactly, the theory is considered exactly solved. This is because, by recursively applying the Operator Product Expansion, the product of two fields inside any correlation function can be systematically replaced by a sum over single fields. Through this contraction procedure, any \(N\)-point correlation function can be reduced to a combination of \((N-1)\)-point functions, and eventually completely decomposed down to known 2-point or 1-point functions.

However, in a generic conformal field theory, these structure constants and scaling dimensions are not known \textit{a priori}. One must rely on sophisticated non-perturbative methods (such as the conformal bootstrap mentioned in the previous chapter) to constrain and compute these numbers.

Furthermore, the physical importance of the OPE is most evident in the short-distance limit, when two fields are brought arbitrarily close to each other (\(x_1 \to x_2\)). In this limit, the expansion typically exhibits a divergence dictated by the power-law denominator \(|x_{12}|^{d_{p}+d_{q}-d_{r}}\). While the OPE naturally contains both regular and singular terms, it is precisely this \textit{singular part} that encodes the most vital algebraic information about the correlation functions and the symmetries of the theory. To fully appreciate how these short-distance singularities dictate the entire structure of the theory, it is highly advantageous to restrict our attention to the incredibly rich case of two dimensions.

\subsection{The Conformal Group in Two Dimensions}

If the dimension of spacetime is larger than two (\(D>2\)), the conformal group is a finite-dimensional Lie group. While it is larger than the standard Poincar\'e group and therefore provides more constraints on the theory, the number of these constraints remains strictly finite. However, in exactly two dimensions, the situation changes drastically and the conformal symmetry becomes infinite-dimensional.

To see this, let us work in 2D Euclidean space (if we start from a Minkowski metric, we can always perform a Wick rotation). In a 2D Minkowski spacetime, the natural variables would be the \textit{lightcone coordinates}; in Euclidean space, the analogous operation is to map the 2D Cartesian plane into the complex plane.
We take a new basis by changing coordinates from the real Cartesian coordinates \(x^1\) and \(x^2\) to the complex coordinates \(z\) and \(\overline{z}\):
\[
    \begin{dcases}
        z = x^1 + i x^2, \\
        \overline{z} = x^1 - i x^2,
    \end{dcases} \implies \begin{dcases}
        x^1 = \frac{z + \overline{z}}{2}, \\
        x^2 = \frac{z - \overline{z}}{2i}.
    \end{dcases}
\]
While physically one coordinate is the complex conjugate of the other, it is mathematically very convenient to treat \(z\) and \(\overline{z}\) as two formally independent variables characterizing a point in \(\mathbb{C}^2\).

In these new coordinates, the flat Euclidean line element becomes:
\[
    \mathrm{d} s^2 = (\mathrm{d} x^1)^2 + (\mathrm{d} x^2)^2 = \mathrm{d} z \mathrm{d} \overline{z},
\]
where \(\mathrm{d} z \mathrm{d} \overline{z}\) also acts as the area element (\(\mathrm{d}^2 x\)) up to a constant factor. The corresponding derivatives with respect to the complex coordinates are:
\[
    \begin{dcases}
        \partial_z = \frac{1}{2}(\partial_1 - i \partial_2), \\
        \partial_{\overline{z}} = \frac{1}{2}(\partial_1 + i \partial_2),
    \end{dcases} \implies \partial_z z = 1, \quad \partial_z \overline{z} = 0,
\]
and similarly for \(\partial_{\overline{z}}\). From the line element \(\mathrm{d}s^2 = g_{\mu\nu}\mathrm{d}x^\mu \mathrm{d}x^\nu\), we can read off the metric tensor in the \((z, \overline{z})\) basis, which takes the off-diagonal form:
\[
    g = \begin{pmatrix}
        0     & 1 / 2 \\
        1 / 2 & 0
    \end{pmatrix}.
\]
Consequently, any field in our theory can be redefined as a function of these two independent variables:
\[
    \varphi(x) \to \varphi(z, \, \overline{z}).
\]

Now, let us consider conformal transformations. A coordinate transformation is conformal if it simply rescales the line element by a strictly positive local factor \(\Omega\):
\[
    (\mathrm{d}s^{\prime})^2 =  \Omega(z, \, \overline{z}) \mathrm{d} s^2.
\]
In complex coordinates, we can achieve this by applying independent transformations to \(z\) and \(\overline{z}\). Let us map them to new variables \(w\) and \(\overline{w}\):
\[
    \begin{dcases}
        z \mapsto w = f(z), \\
        \overline{z} \mapsto \overline{w} = \overline{f}(\overline{z}),
    \end{dcases} \quad \implies \quad \begin{dcases}
        \mathrm{d} w = f^{\prime}(z) \mathrm{d} z, \\
        \mathrm{d} \overline{w} = \overline{f}^{\prime}(\overline{z}) \mathrm{d} \overline{z}.
    \end{dcases}
\]
Here, \(f(z)\) is an analytic (holomorphic) function, meaning its derivative \(f'(z)\) exists, and \(\overline{f}(\overline{z})\) is an anti-holomorphic function. The new line element becomes:
\[
    (\mathrm{d} s^{\prime} )^2 = \mathrm{d} w \mathrm{d} \overline{w} = f^{\prime}(z) \overline{f}^{\prime}(\overline{z}) \mathrm{d} z \mathrm{d} \overline{z},
\]
which perfectly matches the conformal condition with the scale factor completely factorized into a holomorphic and an anti-holomorphic part:
\[
    \Omega(z, \, \overline{z}) = f^{\prime}(z) \overline{f}^{\prime}(\overline{z}).
\]
This reveals a profound result: \textit{any} analytic transformation of the complex plane is a conformal transformation. Geometrically, analytic functions preserve local angles, maintaining the shape of infinitesimally small figures, even though globally they can stretch and severely distort areas (a classic example is the Mercator projection of a sphere, where the North Pole is stretched into an entire infinite line).

Since there are infinitely many analytic functions, the conformal group in two dimensions is \textbf{infinite-dimensional}. This implies we have infinitely many symmetries and, consequently, infinitely many constraints on the algebra of local fields. When a theory is subjected to infinitely many constraints, it often becomes exactly solvable, allowing us to compute quantities far beyond standard perturbation theory.

To rigorously link this back to the conformal Killing equation, consider an infinitesimal conformal transformation:
\[
    \begin{aligned}
        z            & \mapsto z + \epsilon(z),                                  \\
        \overline{z} & \mapsto \overline{z} + \overline{\epsilon}(\overline{z}),
    \end{aligned}
\]
where we define the displacement vector components as:
\[
    \epsilon(z) = \epsilon_1 + i \epsilon_2, \quad \overline{\epsilon}(\overline{z}) = \epsilon_1 - i \epsilon_2.
\]
By definition, an infinitesimal transformation is conformal if it satisfies the conformal Killing equation. In \(D=2\) dimensions, this equation reads:
\[
    \partial_\mu \epsilon_\nu + \partial_\nu \epsilon_\mu = \frac{2}{2} (\partial \cdot \epsilon) \delta_{\mu\nu} = (\partial_1\epsilon_1 + \partial_2\epsilon_2) \delta_{\mu\nu}.
\]
Evaluating this tensor equation component by component (\(\mu=1, \nu=1\) and \(\mu=1, \nu=2\)), we obtain exactly the celebrated \textbf{Cauchy-Riemann equations}:
\[
    \begin{dcases}
        \partial_1 \epsilon_1 = \partial_2 \epsilon_2, \\
        \partial_1 \epsilon_2 = -\partial_2 \epsilon_1.
    \end{dcases}
\]
These equations are the necessary and sufficient condition for \(\epsilon(x^1, x^2) = \epsilon_1 + i \epsilon_2\) to be a complex analytic function \(\epsilon(z)\). This confirms that in two dimensions, the set of all conformal transformations coincides precisely with the set of all holomorphic and anti-holomorphic mappings.

    [...]

\[
    \epsilon(z) = \sum_{n = -\infty}^{\infty} \epsilon_n z^{n + 1},
\]
unless you have a constant function, if it is analytic there must be somewhere in the domain of analiticity in which the function is greater or higher than some value, and outside the domain of analiticity the function can diverge. This laurent series is the most general form of an analytic function, and it is convergent in the annulus defined by the inner and outer radius of convergence.

We can consider \(\epsilon_n(z)\) our basis for the expansion
\[
    \epsilon_n(z) = - \epsilon z^{n + 1},
\]
then the generators of the transformations on scalar fields \(\varphi(x)\) as
\[
    \hat{\ell}_n = - z^{n + 1} \partial_z,
\]
such that
\begin{enumerate}
    \item \(\hat{\ell}_0 = -z \partial_z\) is the generator of \textbf{dilatations},
    \item \(\hat{\ell}_{-1} = -\partial_z\) is the generator of \textbf{translations},
    \item \(\hat{\ell}_1 = -z^2 \partial_z\) is the generator of \textbf{special conformal transformations}.
\end{enumerate}

The dilatation or translation is in \(z\), but we need the \(\overline{z}\) direction, such that we need to define
\[
    \overline{\ell}_n = - \overline{z}^{n + 1} \partial_{\overline{z}}.
\]
So we have
\begin{enumerate}
    \item \(\overline{\ell}_0 = -\overline{z} \partial_{\overline{z}}\) is the generator of \textbf{dilatations},
    \item \(\overline{\ell}_{-1} = -\partial_{\overline{z}}\) is the generator of \textbf{translations},
    \item \(\overline{\ell}_1 = -\overline{z}^2 \partial_{\overline{z}}\) is the generator of \textbf{special conformal transformations}.
\end{enumerate}

Thus if we were to write the four momentum operator in the complex coordinates, we would have
\[
    P = \begin{pmatrix}
        E \\
        p
    \end{pmatrix} = \begin{pmatrix}
        \ell_{-1} + \overline{\ell}_{-1} \\
        i(\ell_{-1} - \overline{\ell}_{-1})
    \end{pmatrix},
\]
where \(E\) is the energy and \(p\) is the momentum. The dilatation operator would be
\[
    D = \hat{\ell}_0 + \overline{\ell}_0,
\]

[...]

If you reason in higher dimensions, the special conformal group is six-dimensional, since
\[
    D=2 \implies \frac{(D+1)(D+2)}{2} = 6.
\]

Imagine to apply the commutator \(\left[ \ell_n, \, \ell_m \right]\) on a test function \(f(z)\), we can find
\[
    \left[ \ell_n, \, \ell_m \right] = \left( (n - m) z^{n + m + 1} \partial_z \right) = (n - m) \ell_{n + m},
\]
which is the Witt algebra, the algebra of the generators of conformal transformations in two dimensions. The same holds for the barred generators \(\overline{\ell}_n\). The Witt algebra is an infinite-dimensional Lie algebra, since \(n\) and \(m\) can take any integer value and be applied on any test function \(f(z)\) in the domain of analyticity. The Witt algebra is the classical version of the Virasoro algebra, which is the quantum version of the conformal algebra in two dimensions. Visaroro was a young phd in Turin when he discovered the algebra, after is supervisor (which already had found the Mobius transformations we are going to introduce later) told him to find the algebra of the generators of conformal transformations in two dimensions, and he found this infinite-dimensional algebra.

This generators, with \(n = -1,0,1\) (barred and not), are six generators which form the algebra of the \(2\times 2\) special linear group \(\mathfrak{sl}(2, \, \mathbb{C})\), which is the double cover of the Lorentz group \(SO(3,1)\). This is the conformal group in two dimensions. It is isomorphic to \(SU(2)\oplus SU(2)\).

We have an infinite number of conformal transformations, since they are all the analytic functions, so how can it be six dimensional? We can map the coordintes of the plane to \(z=e^{w}\), which inverse gives the logarithm of \(z\): we are mapping the plane to a strip (given by the riemann surface of the logarithm), and with periodic boundary conditions we get a cylinder. The conformal transformations on the plane are mapped to conformal transformations on the cylinder, and the six generators of the conformal group are mapped to the six generators of the conformal group on the cylinder. very important (it is a way of compactifying the plane).

Now we can ask which are the transformations starting from a plane and ending in a plane, i.e. which are the transformations that do not change the point at infinity. These are the Mobius transformations, which are given by
\[
    w = \frac{az + b}{cz + d}, \quad ad - bc = 1,
\]
which has six parameters. Thus global conformal transformations are given by the Mobius transformations, which are a subgroup of the infinite-dimensional group of local conformal transformations.

    [...]

We can now infer someting about the fields: after the remapping
\[
    \begin{dcases}
        z \\
        \overline{z}
    \end{dcases} \mapsto \begin{dcases}
        w = f(z) \\
        \overline{w} = \overline{f}(\overline{z})
    \end{dcases},
\]
we can find fields that transforms as
\[
    \dots
\]
which are called \textbf{primary fields}, and they are the building blocks of the theory. We can also understand quasi-primary fields, which are fields that transforms as primary fields under the global conformal transformations (Mobius transformations), but not under the local conformal transformations. There are also other fields which do not have tensorial stucture under the conformal transformations, but they are not important for us. Some of the secondary fields happens to be quasi-primary, but not all of them. The primary fields are the most important fields in the theory and we will distiguish them from the secondary fields, which are obtained by applying the generators of the conformal transformations on the primary fields.

For example, the energy-momentum tensor \(T(z)\) is secondary field, and it would seem lika a bad news since it is the most important field in the theory, but we will see the consequences.

If we compute tge derivatives of \(w\) and \(\overline{w}\) with respect to \(z\) and \(\overline{z}\), we can find
\[
    \dots
\]
and now expanding
\[
    \phi(z, \, \overline{z}) = \phi(w - \epsilon, \, \overline{w} - \overline{\epsilon}) = \phi(w, \, \overline{w}) - \epsilon \partial_w \phi(w, \, \overline{w}) - \overline{\epsilon} \partial_{\overline{w}} \phi(w, \, \overline{w}) + \dots
\]
such that we can find the variation of the field \(\phi\) under the conformal transformation
\[
    \delta \phi = - \epsilon \partial_z \phi - \overline{\epsilon} \partial_{\overline{z}} \phi - h \partial_z \epsilon \phi - \overline{h} \partial_{\overline{z}} \overline{\epsilon} \phi,
\]
where we can see how the two variables are completely decoupled, and the transformation of the field is given by the sum of the transformation in \(z\) and the transformation in \(\overline{z}\).

dilatations\dots
\[
    D = \ell_0 + \overline{\ell}_0 = h + \overline{h} = \Delta,
\]

rotations\dots
\[
    L = \ell_0 - \overline{\ell}_0 = h - \overline{h} = s,
\]

such that we can characterize the fields by their scaling dimension \(\Delta\) and their spin \(s\) (or equivalently by \(h\) and \(\overline{h}\)). A scalar field is spinless, so \(s=0\) and \(h = \overline{h}\). A chiral field is a field that depends only on \(z\) or only on \(\overline{z}\), so it has either \(h=0\) or \(\overline{h}=0\).